\section{Discussion}
\label{sec:discussion}

\paragraph{Binary constraints as a regularizer.}
The central finding reframes binarization. On an encoder without retrieval
pre-training, a sign constraint is not only a compression step but a regularizer
that de-anisotropizes the embedding space. It spreads variance across more axes,
de-saturates the similarity distribution, and improves fine-grained ranking. The
constraint helps most where the float32 space is most anisotropic, which is why
the benefit is ordered RoBERTa $>$ MPNet $>$ BGE and is absent for the
retrieval-pre-trained encoder. The effective dimension is thus a cheap
\emph{diagnostic}: measuring $d_{\mathrm{eff}}$ of an encoder's float32
embeddings, which costs one eigendecomposition and no training, predicts whether
training-time binarization will help or hurt that backbone.

\paragraph{A method-selection guide.}
Because no method dominates, the practical contribution is a rule indexed by two
variables: whether the encoder is pre-trained for retrieval, and the target
memory budget (Table~\ref{tab:cheatsheet}).

\begin{table}[t]
  \centering
  \caption{Recommended compression method by encoder type and compression
  budget, from the Pareto frontiers in Section~\ref{subsec:pareto}.}
  \label{tab:cheatsheet}
  \begin{tabular}{lll}
    \toprule
    \textbf{Budget} & \textbf{Retrieval-pre-trained (BGE)} & \textbf{General-purpose (RoBERTa, MPNet)} \\
    \midrule
    $\leq 8\times$        & INT8 or TurboQuant-4bit (near-lossless) & INT8 or TurboQuant-4bit (near-lossless) \\
    $32\times$            & Post-hoc PQ ($n{=}8$)                   & Training-time BAT (atanh or STE) \\
    $\geq 96\times$       & Stacked MRL$+$PQ                        & Stacked MRL$+$PQ \\
    \bottomrule
  \end{tabular}
\end{table}

\paragraph{Threats to validity.}
We state the limitations plainly. First, evaluation is on NanoBEIR, whose
subsets are small; the backbone-dependent conclusions should be confirmed on
full BEIR~\citep{thakur2021beir} or MTEB~\citep{muennighoff2022mteb} at larger
corpus scale, where compression behavior can shift. Second, each model is
fine-tuned for a single epoch; although the training-trajectory analysis
(Section~\ref{subsec:effdim}) rules out under-training as the cause of the BAT
gain, longer schedules could change absolute scores. Third, we report
compression ratios rather than measured query latency and on-disk index size;
these would close the loop from memory savings to end-to-end cost. Finally, the
effective-dimension account is correlational across five models; a controlled
intervention, such as injecting anisotropy into BGE and observing whether BAT's
relative benefit rises, would strengthen the causal claim.
